\documentclass[conference]{IEEEtran}
\makeatletter
\def\ps@headings{%
\def\@oddhead{\mbox{}\scriptsize\rightmark \hfil \thepage}%
\def\@evenhead{\scriptsize\thepage \hfil \leftmark\mbox{}}%
\def\@oddfoot{}%
\def\@evenfoot{}}
\makeatother
\pagestyle{empty}
\usepackage{url}
\usepackage{graphicx}
\usepackage{amsmath}
\usepackage{amssymb}
\usepackage{multirow}
\usepackage{booktabs}
\usepackage{caption}
\usepackage{hyperref}

\hyphenation{op-tical net-works semi-conduc-tor}

\begin{document}

\title{Noise-Robust Bird Vocalization Classification for\\
Low-Resource Scenarios: A Multi-Paradigm\\
Comparative Study}

\author{
    \IEEEauthorblockN{1\textsuperscript{st} Jianan Zhang}
    \IEEEauthorblockA{\textit{Department of Computer Science} \\
    \textit{Stevens Institute of Technology}\\
    Beijing, China \\
    email: jzhang112@stevens.edu}
    \and
    \IEEEauthorblockN{2\textsuperscript{nd} Jincheng Chen}
    \IEEEauthorblockA{\textit{Department of Computer Science} \\
    \textit{Stevens Institute of Technology}\\
    Guangzhou, China \\
    email: jchen174@stevens.edu}
    \and
    \IEEEauthorblockN{3\textsuperscript{rd} Wenjuan Huang}
    \IEEEauthorblockA{\textit{Department of Computer Science} \\
    \textit{Stevens Institute of Technology}\\
    Zhejiang, China \\
    email: whuang32@stevens.edu}
}

\maketitle

\begin{abstract}
Automated bird vocalization classification from passive acoustic monitoring is critical for biodiversity conservation, yet field recordings present two compounding challenges: environmental noise that masks subtle species-specific features, and extreme long-tail distributions where rare species have minimal training samples. Existing solutions such as BirdNET and Google Perch lack standardized noise robustness evaluation and do not provide cross-paradigm comparisons incorporating traditional machine learning baselines. This work presents a systematic multi-paradigm comparative study on a unified dataset of 129,834 recordings spanning 1,126 bird species from BirdCLEF 2021--2026. We evaluate three approaches: (1) frozen YAMNet with a Dense head on 3072-dimensional clip-level pooled embeddings (mean, max, and std pooling over 1024-dimensional frame-level YAMNet outputs every 0.48 s; encoder not fine-tuned); (2) LightGBM on the same 3072-dimensional pooled features---a cascade decoupling neural feature extraction from tree-based classification, achieving faster inference (7.7 ms/s) than BirdNET ($\sim$10.0 ms/s) with modular flexibility; and (3) FastAI ConvNeXt-Small on mel-spectrograms. All share identical 5-fold CV splits, the same noise-evaluation protocol, CSV \texttt{loss\_class\_weight}, and targeted audio augmentation for low-resource species ($<$15 samples). Noise robustness is evaluated under Gaussian white noise at 5 dB, 0 dB, and $-$5 dB SNR. Our primary contribution is a reproducible evaluation framework with common data splits and noise testing for direct quantitative comparison across heterogeneous ML paradigms---a capability absent from existing bioacoustic tools.
\end{abstract}


\section{Introduction}

Automated bird vocalization classification is essential for large-scale biodiversity monitoring, yet field recordings pose dual challenges: environmental noise masks discriminative spectral features, and species distributions follow an extreme long-tail pattern where rare species have minimal training data [1]--[3]. These compounding factors severely degrade classifier performance, particularly for the very species most in need of conservation monitoring.

Existing off-the-shelf solutions such as BirdNET [1] and Google Perch [2] demonstrate significant limitations under these conditions. BirdNET's accuracy degrades markedly when SNR falls below 3 dB [3], and its recall varies drastically across regions—from PR AUC 0.16--0.23 in North America and Europe to 0.03--0.04 in Africa and Asia [3]. Perch achieves only $\sim$60\% accuracy on BirdCLEF+ 2025 data containing species outside its training set [2]. Critically, neither provides standardized, quantifiable noise robustness evaluation nor cross-paradigm comparisons that include traditional ML baselines.

This project addresses these gaps through three pipelines on a unified 129,834-sample dataset with controlled noise evaluation: (1) frozen YAMNet (MobileNetV1, AudioSet [4]) with a Dense head on 3072-dimensional clip-level pooled embeddings (aggregated from 1024-dimensional frame-level outputs); (2) LightGBM on the same 3072-dimensional pooled features---a two-stage cascade decoupling neural feature extraction from tree-based classification, enabling GPU-free training and modular experimentation; and (3) FastAI ConvNeXt-Small on mel-spectrograms. All share identical 5-fold CV, the same SNR evaluation protocol, CSV \texttt{loss\_class\_weight}, and targeted audio augmentation for species with fewer than 15 training samples. Evaluation uses clean audio plus Gaussian noise at 5 dB, 0 dB, and $-$5 dB SNR. While we use Gaussian white noise as a controlled, reproducible baseline, qualitatively different field noise (wind, rain, insect interference) is discussed in Section V.

Our contributions are threefold: (i) a reproducible cross-paradigm evaluation framework with common data splits and standardized noise testing; (ii) empirical insights into trade-offs between frozen embedding classifiers, embedding-based gradient boosting, and spectrogram-based classification---with emphasis on the YAMNet+LightGBM cascade as a computationally efficient alternative to monolithic systems; and (iii) a practical framework combining shared CSV-based class weighting and targeted low-resource augmentation with standardized evaluation for low-resource bioacoustic classification.


\section{Related Work}

\subsection{BirdNET and Google Perch: Off-the-Shelf Solutions}

BirdNET [1] (Cornell Lab/Chemnitz University) uses an EfficientNetB0-like backbone with 1024-dimensional embeddings, classifying 48 kHz audio into 3,337 classes [7]. Despite its widespread adoption, it exhibits: (1) severe performance degradation below 3 dB SNR [3]; (2) highly heterogeneous recall across regions (0.24--0.52) and species rarity levels [3]; and (3) label noise and dataset imbalances as primary error sources [8]. Google Perch [2] (EfficientNet-B1) achieves state-of-the-art on BirdSET and BEANS benchmarks [9], [10] but only $\sim$60\% accuracy on BirdCLEF+ 2025 data with out-of-training-set species. Critically, neither provides standardized controllable noise testing pipelines.

\subsection{BirdCLEF Competitions: Persistent Challenges}

The BirdCLEF series (2021--2026) consistently benchmarks core challenges: long-tailed distributions, weak labeling (recording-level vs. segment-level), and domain shift between Xeno-Canto training and PAM test soundscapes [5], [11]. BirdCLEF+ 2025 saw foundation models like BEATs achieve 0.913 F1 with pseudo-labeling [12]. These competitions underscore that no integrated solution spans data, model, and evaluation dimensions for low-resource noisy environments.

\subsection{Noise Robustness and Low-Resource Classification}

Environmental noise remains a formidable obstacle. Noise reduction preprocessing can help but occasionally degrades performance by removing informative features [13]. Data augmentation with real environmental noise under varying SNR has proven effective [5], [14], and generative approaches (ACGANs, DDPMs) show promise [15]. For low-resource species, transfer learning is essential: avian-specific embeddings outperform general audio embeddings [16], and AudioSet pre-training followed by fine-tuning boosts low-resource performance [17], [16]. Few-shot approaches (Siamese networks) target rare species absent from BirdNET/Perch [18]. However, existing studies rarely provide systematic cross-paradigm comparisons on identical datasets.

\subsection{YAMNet and Audio Transfer Learning}

YAMNet [4], [17] (MobileNetV1, AudioSet-trained) outputs 1024-dimensional frame-level embeddings at 0.48-second intervals. In our pipelines, these frames are aggregated by mean, maximum, and standard deviation into 3072-dimensional clip-level pooled features for downstream classifiers. It has been adapted for speech recognition, urban noise, and industrial monitoring [19]. While domain-specific models (Perch, BirdNET) outperform general-purpose ones on complex calls [16], [20], YAMNet provides a computationally accessible baseline for raw audio classification.

\subsection{Research Gap Summary}

The literature reveals six critical gaps: (1) no unified framework comparing tabular ML, image CNN, and raw-audio transfer learning on identical datasets with standardized noise evaluation; (2) off-the-shelf tools lack quantifiable controllable noise pipelines; (3) no systematic cross-paradigm comparison for rare species; (4) no reproducible controlled-variable experiment separating clean baseline from gradient noise; (5) BirdCLEF solutions rarely integrate data, model, and evaluation dimensions; and (6) no systematic evaluation of hybrid pipelines combining deep embeddings with lightweight tree classifiers—which can outperform monolithic systems while being faster and more deployable. Our YAMNet+LightGBM cascade directly addresses gap (6), achieving $\sim$7.7 ms/s embedding latency (faster than BirdNET's $\sim$10.0 ms/s) with modular flexibility.


\section{Our Solution}

\subsection{Dataset}

\subsubsection{Data Source and Collection}

Compiled from BirdCLEF 2021--2026, aggregating Xeno-Canto and iNaturalist recordings [5], [6]. Raw merge of six annual CSVs: 183,239 records. Field-level deduplication (zero-information-loss merge) yielded 167,308 unique recordings.

\subsubsection{Data Cleaning and Quality Control}

Two filters applied:
\begin{enumerate}
    \item Remove numeric-only \texttt{primary\_label} (non-avian taxa: Amphibia, Insecta, Mammalia, Reptilia)—excludes 101 non-bird species;
    \item Discard rating $\leq$ 0.5 (unrated/extremely low-quality).
\end{enumerate}
Final cleaned population: \textbf{144,250 recordings across 1,126 bird species}.

SNR proxy tier derived from rating: High ($\geq$4.0), Medium (2.0--3.9), Heavy ($<$2.0). Split: 90/10 with SNR stratification (seed=42). Species with $<$5 recordings retained entirely in training. Tables I--II summarize the split.

\begin{table}[!t]
\centering
\caption{Final Dataset Split (V2.0)}
\label{tab:split}
\begin{tabular*}{\columnwidth}{@{\extracolsep{\fill}}lrrl}
\toprule
\textbf{Split} & \textbf{Records} & \textbf{Species} & \textbf{Use} \\
\midrule
Cleaned population & 144,250 & 1,126 & Source \\
Training pool & 129,834 & 1,126 & Development \& CV \\
CV validation (per fold) & $\sim$25,967 & 1,126 & In-fold validation (5-fold SNR-stratified CV) \\
Fixed holdout test & 14,416 & 1,093 & Final evaluation \\
\bottomrule
\end{tabular*}

\vspace{2\baselineskip}

\captionof{table}{SNR Tier Distribution Across\\ Train/Test Splits}
\label{tab:snr}
\begin{tabular}{lrr}
\toprule
\textbf{SNR Tier} & \textbf{Training} & \textbf{Test} \\
\midrule
High ($\geq$4.0) & 86,731 & 9,631 \\
Medium (2.0--3.9) & 41,011 & 4,553 \\
Heavy ($<$2.0) & 2,092 & 232 \\
\bottomrule
\end{tabular}
\end{table}

5-fold SNR-stratified CV: each fold $\sim$103,853 training, $\sim$25,967 validation. No filename overlap between train/test or train/val within folds. This represents a \textbf{10$\times$ scale-up} from the original proposal (12,400 $\rightarrow$ 129,834), enabling robust long-tail evaluation.

\subsubsection{Long-Tail Distribution and Mitigation}

Species distribution: median $\sim$72 training samples/species, minimum 1, head class 1,300. \textbf{Shared long-tail mitigation (all three pipelines)}:

\begin{enumerate}
    \item \textbf{Balanced sampling}: $w_{\text{sampler}} = 1 / N_c$, ensuring equal per-species batch probability (exported for LightGBM; not used in Dense-head training).
    
    \item \textbf{Mild class-weighted loss}:
    \begin{equation}
        w_{\text{loss}}(c) = \operatorname{clip}\left(\frac{1/\sqrt{N_c}}{\mu_w},\, 0.25,\, 4.00\right)
        \label{eq:loss_weight}
    \end{equation}
    Inverse-square-root is milder than $1/N_c$; clipping ensures stability. Species with fewer than 15 training samples are treated as low-resource tail classes, consistent with BirdCLEF findings that such species exhibit near-random performance [11].
    
    \item \textbf{Targeted audio augmentation}: 159 species (14.1\%) with $<$15 training examples receive four on-the-fly waveform transformations implemented in \texttt{audio\_augmentation.py}: time stretching (0.85$\times$--1.10$\times$), pitch shifting ($\pm$1--$\pm$2 semitones), additive Gaussian noise (5--15 dB SNR), and volume scaling (0.70$\times$--1.30$\times$), capped at 16 variants per recording and 50 new variants per species, with a per-species target of 15 samples. This augmentation is enabled in all three reported pipelines: YAMNet applies \texttt{expand\_with\_augmentation} before embedding extraction and Dense-head training; LightGBM trains on 3072-dimensional pooled embeddings computed from the same augmented training folds; FastAI generates augmented mel-spectrogram PNGs via \texttt{augmentation\_glue.py} before ConvNeXt training.
\end{enumerate}

\subsubsection{Evaluation Protocol}

Fixed test set (14,416) reserved for final evaluation only. Metrics:
\begin{enumerate}
    \item clean Top-1 accuracy;
    \item noise robustness at 5 dB, 0 dB, $-$5 dB SNR (Gaussian white noise, deterministic per-sample seeding with base seed=42, peak-normalized);
    \item per-SNR-tier stratified accuracy;
    \item inference latency (ms/sample) and peak GPU memory.
\end{enumerate}

\subsection{Machine Learning Algorithms}

\subsubsection{Frozen YAMNet Embeddings with Dense Head}

\textbf{Rationale.} MobileNetV1 pre-trained on AudioSet's 521 categories [4] provides general-purpose acoustic features transferable to bird vocalizations. AudioSet includes environmental noise [17], giving inherent noise-robust representations. We freeze the encoder and train only a lightweight head on pooled embeddings [16].

\textbf{Design.} 16 kHz mono, 5-second clips (center-crop or zero-pad), peak-normalized. (1) Frozen YAMNet encoder (\texttt{trainable=False}) outputs 1024-dimensional frame-level embeddings every 0.48 s; mean, max, and std pooling across frames $\rightarrow$ 3072-dimensional clip-level feature vector. (2) Dense head: Dense(256, ReLU) $\rightarrow$ Dropout(0.3) $\rightarrow$ Dense(1126, Softmax). Adam ($\text{lr}=10^{-3}$), up to 50 epochs, early stopping (patience=8), ReduceLROnPlateau (factor=0.5, patience=4), \texttt{loss\_class\_weight} as training sample weights. Low-resource species ($<$15 samples) are expanded via \texttt{expand\_with\_augmentation} before embedding extraction. Embeddings are cached once and exported per fold for LightGBM. Noise evaluation re-computes clip-level pooled embeddings on noisy test audio.

\subsubsection{LightGBM: Gradient-Boosted Decision Trees on YAMNet Embeddings}

\textbf{Rationale.} Gradient-boosted trees complement deep learning: no GPU required, fast CPU training, native feature importance, robust to hyperparameters. However, trees cannot operate on raw audio. Our \textbf{cascade architecture} bridges this: YAMNet's frozen encoder extracts embeddings; LightGBM classifies. This decoupling offers three advantages: (i) embedding extraction runs once per file; (ii) classification is GPU-free, trainable in minutes with class-weighted loss; (iii) modular—same embeddings reusable for different classifiers.

\textbf{Audio Preprocessing.} 16 kHz mono OGG loaded directly via librosa (no WAV conversion). Fixed 5.0 seconds (zero-pad if shorter; center-crop if longer), peak-normalized to [$-1$, 1]. \textbf{No noise reduction}—evaluation requires measuring robustness against controlled additive noise.

\textbf{Feature Extraction.} The frozen YAMNet encoder produces 1024-dimensional frame-level embeddings at 0.48 s intervals ($\sim$10 frames per 5 s clip). These are pooled by \textbf{mean, maximum, and standard deviation} into a \textbf{3072-dimensional clip-level feature vector}---not the raw 1024-d frame outputs. This captures both average spectral profile and temporal variability—critical for dynamic vocalization patterns. Compared to handcrafted features (MFCC statistics, spectral centroid), YAMNet pooled features preserve deeper semantic information; compared to flattened spectrograms ($\sim$64,000 dims), 3072 dims makes tree learning tractable.

\textbf{Model Design.} LightGBM: 1,500 rounds with early stopping (patience=50), learning rate 0.05, max\_depth=6, num\_leaves=31, L1/L2 regularization ($\alpha=0.1$, $\lambda=0.1$), \texttt{loss\_class\_weight} as sample weights. CPU inference only. Outputs 1,126-dimensional raw class scores under the multiclass objective (margin logits, not native softmax). For probability-based metrics, scores are normalized via softmax at evaluation time.

\textbf{Methodological value:} Establishes performance floor for CPU-only tools with high-quality embeddings---quantifying accuracy-per-compute trade-offs against deep neural classifiers on the same embeddings.

\subsubsection{FastAI ConvNeXt: Mel-Spectrogram Image Classification}

\textbf{Rationale.} Mel-spectrograms expose harmonic structures, formant patterns, and temporal modulations. ConvNeXt-Small provides strong performance with ImageNet transfer learning [14], [11]. Long-tail imbalance is handled with the same CSV \texttt{loss\_class\_weight} as YAMNet/LightGBM, supplemented by spectrogram-level MixUp and RandomErasing. Low-resource waveform augmentation is applied upstream via \texttt{augmentation\_glue.py} before PNG caching.

\textbf{Design.} 16 kHz mono; highest-RMS 5 s segment; 128 mel bands (\texttt{n\_fft=1024}, hop 512, 500--8000 Hz) $\rightarrow$ Viridis RGB images at 224$\times$224 (cached offline, including augmented low-resource samples). ConvNeXt-Small, FP16, CSV \texttt{loss\_class\_weight} as sample weights, MixUp ($\alpha=0.4$), RandomErasing ($p=0.4$). \texttt{fine\_tune}: 3 frozen epochs + 12 full-network epochs (\texttt{base\_lr=$2\times10^{-3}$}). Metrics: Top-1 and macro-F1.

\subsection{Implementation Details}

All pipelines share:

\begin{itemize}
    \item Identical 5-fold CV splits (SNR-stratified)
    \item Identical noise evaluation: Gaussian white noise at 5/0/$-$5 dB, deterministic per-sample seeding (base seed=42)
    \item Same metrics: Top-1, Top-5, macro-F1, balanced accuracy, stratified by SNR tier and class frequency group
    \item Inference measurement: 5 warm-up + 50 timed iterations; GPU memory profiling for FastAI fine-tuning (YAMNet embedding extraction uses $<$1\,GB as a one-time cost); CPU-only for LightGBM
    \item Shared long-tail controls: all pipelines use CSV \texttt{loss\_class\_weight} and targeted audio augmentation for species with $<$15 samples; FastAI additionally applies MixUp and RandomErasing at the spectrogram level
    \item Embedding terminology: YAMNet outputs 1024-d frame-level embeddings; Dense/LightGBM consume 3072-d clip-level pooled features (mean/max/std)
\end{itemize}


\section{Comparison}

\subsection{Experimental Results}

[To be completed after model training and evaluation. Will present: (1) comparative Top-1, Top-5, macro-F1 on clean test audio; (2) noise robustness degradation curves (5 dB, 0 dB, $-$5 dB); (3) per-SNR-tier and per-class-frequency-group stratified performance; (4) inference latency and GPU memory benchmarks.]

\subsection{Comparison with Existing Solutions}

[To be completed. Will compare best-performing model against BirdNET and Perch baselines on the same test set, analyzing trade-offs between transfer learning, YAMNet embedding-based gradient boosting, and CNN image classification.]


\section{Future Directions}

\begin{enumerate}
    \item \textbf{Multi-label classification}:\linebreak[0] Incorporate~\texttt{secondary\_\allowbreak labels} for partially labeled co-occurrences, reducing label noise and improving tail-class performance.
    \item \textbf{Attention-based embedding aggregation}: Replace simple mean/max/std pooling with learnable attention to preserve temporally localized discriminative patterns.
    \item \textbf{Extended noise conditions}: Evaluate against wind, rain, and insect noise—qualitatively different from Gaussian.
    \item \textbf{Model distillation}: Distill YAMNet+LightGBM cascade into a single lightweight edge-deployable model.
    \item \textbf{Active learning for rare species}: Identify most informative unlabeled recordings for human annotation, maximizing labeling efficiency for species with $<$5 samples.
\end{enumerate}


\section{Conclusion}

[To be completed after experimental results are obtained. This section will summarize: whether the noise-robust bird vocalization classification problem is adequately addressed by our multi-paradigm approach; which algorithm(s) demonstrate superior performance under specific conditions (clean vs. noisy, head vs. tail classes); which shared techniques (CSV \texttt{loss\_class\_weight}, targeted low-resource augmentation, frozen AudioSet embeddings vs. ConvNeXt spectrogram modeling) contribute most to performance gains; and broader implications for low-resource bioacoustic classification system design.]


\section*{Acknowledgment}

The authors would like to thank their advisor for guidance and the Kaggle community for providing the BirdCLEF datasets.


\begin{thebibliography}{99}

\bibitem{kahl2021birdnet}
S. Kahl, C. M. Wood, M. Eibl, and H. Klinck, ``BirdNET: A deep learning solution for avian diversity monitoring,'' \textit{Ecological Informatics}, vol. 61, p. 101236, 2021.

\bibitem{van2025perch}
B. van Merri\"{e}nboer et al., ``Perch 2.0: The bittern lesson for bioacoustics,'' \textit{arXiv:2508.04665}, 2025.

\bibitem{funosas2026global}
D. Funosas et al., ``A global assessment of BirdNET performance: Differences among continents, biomes, and species,'' \textit{Ecological Indicators}, vol. 182, p. 114550, 2026.

\bibitem{yamnet2024}
S. Hershey et al., ``CNN architectures for large-scale audio classification,'' in \textit{Proc. IEEE ICASSP}, 2017. \textit{(YAMNet implements this architecture on AudioSet [17]; public model release: TensorFlow Hub, 2019.)}

\bibitem{birdclef2026}
Cornell Lab of Ornithology, ``BirdCLEF / BirdCLEF+ competitions (2021--2026),'' Kaggle, 2026.

\bibitem{xenocanto}
Xeno-Canto Foundation, ``Xeno-Canto: Sharing bird sounds from around the world,'' \url{https://xeno-canto.org/}.

\bibitem{kahl2025birdnetv2}
S. Kahl, C. M. Wood, and H. Klinck, ``BirdNET model v2.4,'' Zenodo, 2025.

\bibitem{assarzadeh2025investigating}
M. Assarzadeh, P. Atufi, M. Eibl, and S. Kahl, ``Investigating misclassification in bird sounds: The adverse effect of label noise on BirdNET,'' \textit{Ecological Informatics}, vol. 92, p. 103494, 2025.

\bibitem{rauch2024birdset}
L. Rauch et al., ``BirdSet: A large-scale dataset for audio classification in avian bioacoustics,'' \textit{arXiv:2403.10380}, 2024.

\bibitem{hagiwara2023beans}
M. Hagiwara et al., ``BEANS: The benchmark of animal sounds,'' in \textit{Proc. IEEE ICASSP}, 2023, arXiv:2210.12300.

\bibitem{witting2024addressing}
E. Witting, H. de Heer, J. Lim, C. T. Kopar, and K. S\'{a}ndor, ``Addressing the challenges of domain shift in bird call classification for BirdCLEF 2024,'' in \textit{CEUR Workshop Proc.}, vol. 3740, 2024.

\bibitem{chen2023beats}
S. Chen et al., ``BEATs: Audio pre-training with acoustic tokenizers,'' in \textit{ICML 2023}, arXiv:2212.09058.

\bibitem{denton2022improving}
T. Denton, S. Wisdom, and J. R. Hershey, ``Improving bird classification with unsupervised sound separation,'' in \textit{Proc. IEEE ICASSP}, 2022, arXiv:2110.03209.

\bibitem{miyaguchi2023transfer}
A. Miyaguchi, N. Zhong, M. Gustineli, and C. Hayduk, ``Transfer learning with semi-supervised dataset annotation for birdcall classification,'' in \textit{CLEF 2023 Working Notes}.

\bibitem{gibbons2024generative}
A. Gibbons, E. King, I. Donohue, and A. Parnell, ``Generative AI-based data augmentation for improved bioacoustic classification in noisy environments,'' \textit{arXiv:2412.01530}, 2024.

\bibitem{ghani2023global}
B. Ghani, T. Denton, S. Kahl, and H. Klinck, ``Global birdsong embeddings enable superior transfer learning for bioacoustic classification,'' \textit{Scientific Reports}, vol. 13, p. 22876, 2023.

\bibitem{gemmeke2017audioset}
J. F. Gemmeke et al., ``Audio set: An ontology and human-labeled dataset for audio events,'' in \textit{ICASSP 2017}.

\bibitem{zhong2023comparison}
M. Zhong, J. LeBien, M. Campos-Cerqueira, T. M. Aide, R. Dodhia, and J. Lavista Ferres, ``A comparison of convolutional neural networks and few-shot learning in classifying long-tailed distributed tropical bird songs,'' \textit{bioRxiv}, 2023.

\bibitem{liu2025yamnet}
L. Liu et al., ``YAMNet-based transfer learning for compact noise classification in urban and wireless systems,'' \textit{EURASIP J. Wireless Commun. Netw.}, vol. 2025, no. 1, p. 74, 2025.

\bibitem{vggish2024}
``VGGish: Audio feature embedding model,'' Kaggle Models, 2024.

\end{thebibliography}

\end{document}